\ifdefined\gkver\else\def\gkver{student}\fi
\documentclass[\gkver]{gaokaozhenti}
\begin{document}

\chapter{2026年上海}

\begin{enumerate}
\item
%% number: 1
%% paperName: 2026年上海等级考第 1 题 15 分
%% typeId: 9
%% type: 综合题
%% chapter: 曲线运动
%% point: 人造地球卫星
%% method: 无
%% score: 15
%% degree: 550
%% duplicateId: 0
%% body:
\textbf{空间科技}

近年来，我国在深空探测、卫星应用等领域成就斐然。
\begin{enumerate}
	\item
	空间核反应堆中的典型链式反应之一为 $\ce{^{235}_{92}U} + \ce{^{1}_{0}n} \to \ce{^{141}_{56}Ba} + \ce{^{92}_{36}Kr} +$ \tkanswer[1.5cm]{\ce{3^{1}_{0}n}}。

	\fourchoices[answer=B]
	{$\ce{^{1}_{0}n}$}
	{$3\ce{^{1}_{0}n}$}
	{$\ce{^{3}_{1}H}$}
	{$3\ce{^{0}_{1}e}$}

	\item
	空间核反应堆中某次核反应前的总质量为 $m_{1}$，反应后的总质量变为 $m_{2}$，真空中的光速为 $c$，则该次核反应释放的能量为 \tkanswer[2cm]{$(m_{1}-m_{2})c^{2}$}。

	\item
	一绕地卫星变轨过程中，下降高度 $\Delta h$，引力势能变化量为 $\Delta E_{p1}$，继续下降高度 $\Delta h$，引力势能再次变化，变化量为 $\Delta E_{p2}$，则 \tkanswer[1cm]{A}。

	\threechoices[answer=A]
	{$|\Delta E_{p1}| < |\Delta E_{p2}|$}
	{$|\Delta E_{p1}| = |\Delta E_{p2}|$}
	{$|\Delta E_{p1}| > |\Delta E_{p2}|$}

	\item
	一质量为 $m$ 的卫星绕地球做半径为 $r$ 的匀速圆周运动时，其动能为 \tkanswer[2cm]{$\frac{GMm}{2r}$}。（地球质量为 $M$，引力常量为 $G$）

	\item
	地球赤道上方有一带正电的绕地卫星相对于地面由西向东运动，地磁场对其产生的洛伦兹力 \tkanswer[1cm]{A}。

	\fourchoices[answer=A]
	{向上}
	{向下}
	{向东}
	{向西}
\end{enumerate}
%% answer:
\jdanswer{
\begin{enumerate}
	\item
	B

	\item
	$(m_{1}-m_{2})c^{2}$

	\item
	A

	\item
	$\frac{GMm}{2r}$

	\item
	A
\end{enumerate}
}
%% memo:
\memoanswer{
1．设所求质量数为 $A$、电荷数为 $Z$，根据质量数守恒、电荷数守恒可得：

$A=(235+1)-(141+92)=3$，$Z=(92+0)-(56+36)=0$。

正确选项为 B。

2．根据爱因斯坦质能方程可得 $E=\Delta mc^{2}=(m_{1}-m_{2})c^{2}$。

3．引力势能变化量 $\Delta E_{p}$ 的大小对应万有引力做的功。根据万有引力定律，下降过程中万有引力增大，下降同样高度 $\Delta h$，第二次引力做的正功多，引力势能变化量大，即 $|\Delta E_{p1}|<|\Delta E_{p2}|$。正确选项为 A。

\onepicture[w=5.0cm]{\includesvg{figs/01_an1}}

另解：由引力势能的表达式 $E_{p}=-\frac{GMm}{r}$ 可以绘制出 $E_{p}-r$ 图像，从绘制出的图像可以看出，在 $r$ 减小的过程中，若横坐标变化量 $\Delta h$ 相同，纵坐标的变化量并不相同，且满足 $|\Delta E_{p1}|<|\Delta E_{p2}|$。

4．根据万有引力提供向心力可得 $G\frac{Mm}{r^{2}}=m\frac{v^{2}}{r}$，由此式可解得卫星的动能 $E_{k}=\frac{1}{2}mv^{2}=\frac{GMm}{2r}$。

5．赤道正上方的地磁场由南向北，带正电的卫星自西向东飞行，根据左手定则可得洛伦兹力的方向向上（指向天空）。

\onepicture[w=4.0cm]{\includesvg{figs/01_an2}}

正确选项为 A。
}

\item
%% number: 2
%% paperName: 2026年上海等级考第 2 题 18 分
%% typeId: 9
%% type: 综合题
%% chapter: 电场
%% point: 带电粒子的偏转
%% method: 无
%% score: 18
%% degree: 450
%% duplicateId: 0
%% body:
\textbf{光刻技术}

极紫外光刻是重要的微加工技术之一。
\begin{enumerate}
	\item
	一单色光经过双缝投射到光刻胶涂层上，得到如图所示的干涉图样。为减小条纹间距，可减少 \tkanswer[1.5cm]{AD}（多选）。

	\onepicture[w=4.5cm, align=right]{figs/02a.jpg}

	\fourchoices[answer=AD]
	{光波长}
	{光强}
	{双缝间距}
	{双缝到涂层的距离}

	\item
	产生光刻所需极紫外光的简化原理图如图所示。平行板电容器 $C$ 与电阻 $R$、电源 $E$ 串联，带电锡液滴在重力作用下从 $a$ 点由静止下落，穿过竖直放置的电容器 $C$ 运动到 $b$ 点处被激光照射，产生处于高能级态的锡离子之后，锡离子向低能级态跃迁，辐射出极紫外光。（所有锡液滴视为完全相同，不考虑电容器的边缘效应）

	\onepicture[h=4.3cm, align=right]{\includesvg{figs/02b}}

	\begin{enumerate}
		\item
		若电源输出电压为 $U$，电容器两极板间距为 $d$，则极板间电场的电场强度大小为 \tkanswer[1.6cm]{$\frac{U}{d}$}。

		\item
		设一锡液滴进入电容器时电势能为 $E_{p1}$，其所在位置电势为 $\varphi_{1}$；该液滴即将离开电容器时电势能为 $E_{p2}$，所在位置电势为 $\varphi_{2}$，则 \tkanswer[1cm]{D}。

		\fourchoices[answer=D]
		{$E_{p1} < E_{p2}$，$\varphi_{1} > \varphi_{2}$}
		{$E_{p1} > E_{p2}$，$\varphi_{1} > \varphi_{2}$}
		{$E_{p1} < E_{p2}$，$\varphi_{1} < \varphi_{2}$}
		{$E_{p1} > E_{p2}$，$\varphi_{1} < \varphi_{2}$}

		\item
		若锡滴都落在 $b$ 点右侧，在不改变其他条件的情况下，为使锡液滴能落至 $b$ 点，可将 \tkanswer[1.5cm]{CE}（多选）。

		\sixchoices[answer=CE]
		{电阻 $R$ 的阻值调小}
		{电阻 $R$ 的阻值调大}
		{电容器 $C$ 整体上移}
		{电容器 $C$ 整体下移}
		{电容器 $C$ 右极板向左移}
		{电容器 $C$ 右极板向右移}

		\item
		锡离子从高能级态向低能级态跃迁发出波长为 $\lambda$ 的极紫外光，则这两个能级的能量差为 \tkanswer[1.8cm]{$\frac{hc}{\lambda}$}。（普朗克常量为 $h$，真空中的光速为 $c$）
	\end{enumerate}

	\item
	光刻技术新方案采用氦原子束代替极紫外光。为使氦原子的物质波波长小于极紫外光的波长 $\lambda$，氦原子的速度大小至少为 \tkanswer[1.6cm]{$\frac{h}{m\lambda}$}。（氦原子质量为 $m$，普朗克常量为 $h$）

	\fourchoices[answer=C]
	{$\frac{m}{h\lambda}$}
	{$\frac{h\lambda}{m}$}
	{$\frac{h}{m\lambda}$}
	{$\frac{mh}{\lambda}$}
\end{enumerate}
%% answer:
\jdanswer{
\begin{enumerate}
	\item
	AD

	\item
	（1）$\frac{U}{d}$     （2）D     （3）CE     （4）$\frac{hc}{\lambda}$

	\item
	C
\end{enumerate}
}
%% memo:
\memoanswer{
1．根据双缝干涉条纹间距公式 $\Delta x=\frac{L}{d}\lambda$ 可知，要减小条纹间距 $\Delta x$，可以减小光的波长 $\lambda$ 或减小双缝到涂层的距离 $L$。

正确选项为 AD。

2．（1）匀强电场中电场强度和电势差的关系是 $E=\frac{U}{d}$。

（2）电荷在电场中运动时，电场力做正功，导致电势能减少，可得 $E_{p1} > E_{p2}$；由电路图可知电容器左极板带正电，因此极板间的电场线方向向右，液滴在此区域向左逆电场线方向运动，电势升高，即 $\varphi_{1} < \varphi_{2}$。

正确选项为 D。

（3）电路稳定时两板间电压等于电源电动势，与电阻无关，A、B 选项错误；

液滴在竖直方向做自由落体运动，若电容器整体上移，则在板间运动的时间 $t$ 变大，离开电场的竖直速度 $v_{y}$ 减小；水平方向做匀加速运动，水平距离 $x=\frac{1}{2}a_{x}t^{2}$、水平速度 $v_{x}=a_{x}t$ 也随之变大，与竖直方向的夹角 $\tan\theta=\frac{v_{x}}{v_{y}}$ 也增大，使液滴离开电场后会比之前更偏向左运动，能落至 $b$ 点。选项 C 正确 D 错误。

电容器右极板向左移动，两板间距离减小，导致水平方向加速度 $a_{x}$ 增大，水平距离 $x$、与竖直方向的夹角 $\tan\theta$ 都增大，使液滴离开电场后会比之前更偏向左运动，能落至 $b$ 点。选项 E 正确 F 错误。

正确选项为 CE。

（4）根据玻尔理论的频率条件，两个能级的能量差 $\Delta E=h\nu=\frac{hc}{\lambda}$。

3．根据德布罗意波长公式，氦原子的物质波波长 $\lambda=\frac{h}{mv}$，由 $\frac{h}{mv}<\lambda$ 解得速度 $v>\frac{h}{m\lambda}$。

正确选项为 C。
}

\item
%% number: 3
%% paperName: 2026年上海等级考第 3 题 13 分
%% typeId: 9
%% type: 综合题
%% chapter: 电磁感应
%% point: 法拉第电磁感应定律
%% method: 无
%% score: 13
%% degree: 450
%% duplicateId: 0
%% body:
\textbf{电吉他}

电吉他拾音器基本结构如图所示。磁铁附近的钢质琴弦被磁化，弦的振动会在闭合线圈中产生感应电流。

\onepicture[h=3.5cm, align=right]{\includesvg{figs/03a}}
\begin{enumerate}
	\item
	拨动琴弦发出一 E2 音（$82 \UHz$），该频率的波动从琴弦传至空气时，\tkanswer[1cm]{D}。

	\fourchoices[answer=D]
	{频率改变，波长不变}
	{频率不变，波长不变}
	{频率改变，波长改变}
	{频率不变，波长改变}

	\item
	琴弦上某质点的振动情况如图所示。$t=1.0 \UmsT$ 时，该质点正在向 $x$ 轴 \tkanswer[1cm]{B}。

	\onepicture[w=5.0cm, align=right]{\includesvg{figs/03b}}

	\fourchoices[answer=B]
	{正方向加速运动}
	{负方向加速运动}
	{正方向减速运动}
	{负方向减速运动}

	\item
	电吉他中有一 $LC$ 振荡电路，以图（a）所示电流 $i$ 方向为正方向。某时刻电容器 $C$ 正在充电且上极板带正电，则该时刻可能对应图（b）中的 \tkanswer[1cm]{A}。

	\onepicture[w=7.0cm, align=right]{\includesvg{figs/03c}}

	\fourchoices[answer=A]
	{$t_{1}$}
	{$t_{2}$}
	{$t_{3}$}
	{$t_{4}$}

	\item
	穿过某拾音器线圈的磁通量 $\Phi$ 随时间 $t$ 的变化关系如图所示，线圈中感应电动势变化的频率为 \tkanswer[1.2cm]{$263$} $\UHz$。若该线圈有 500 匝，则感应电动势最大值为 \tkanswer[1.4cm]{$0.0785$} $\UV$。（结果均保留 3 位有效数字，$1\UuWb=1\times10^{-6}\UWb$，$1\UmsT=1\times10^{-3}\Us$）

	\onepicture[w=9.0cm, align=right]{\includesvg{figs/03d}}
\end{enumerate}
%% answer:
\jdanswer{
\begin{enumerate}
	\item
	D

	\item
	B

	\item
	A

	\item
	263，0.0785（误差有一定范围）
\end{enumerate}
}
%% memo:
\memoanswer{
1．频率由波源决定，波从一种介质进入另一种介质，频率不变；而波速取决于介质，琴弦上传播的机械波传入空气中形成的声波，波速会发生改变。

正确选项为 D。

2．由振动图像可知，质点在 $t=1.0 \UmsT$ 时正处在正方向位移处，向平衡位置移动。而质点在平衡位置速度最大，因此该质点正在向 $x$ 轴负方向加速运动。

正确选项为 B。

3．电容器充电时，电流应流入正极板，因此这个时刻回路中电流方向为逆时针方向，即正方向，对应图中的 $t_{1}$ 时刻。

正确选项为 A。

4．本题原卷及材料中未提供第 4 问详解，待补充。
}

\item
%% number: 4
%% paperName: 2026年上海等级考第 4 题 13 分
%% typeId: 6
%% type: 计算题
%% chapter: 光学
%% point: 全反射
%% method: 无
%% score: 13
%% degree: 450
%% duplicateId: 0
%% body:
\textbf{压缩空气引火仪}

如图，取少许干燥的疏松硝化棉，放入压缩空气引火仪的透明圆筒内，以不同速度压下活塞，观察密闭筒内发生的现象。

\onepicture[h=4.5cm, align=right]{\includesvg{figs/04a}}
\begin{enumerate}
	\item
	缓慢压下活塞，筒内气体视为温度不变的理想气体，硝化棉未被点燃。

	\begin{enumerate}
		\item
		将筒内气体的压强、体积、热力学温度分别记为 $p$、$V$、$T$。以下图线可能描述此过程中气体状态量变化关系的有 \tkanswer[1.5cm]{AD}（多选）。

		\onepicture[w=13cm, align=right]{\includesvg{figs/04b}}

		\item
		当筒内气体体积压缩为初始体积的一半时，与初始状态相比，筒内气体的以下物理量发生变化的是 \tkanswer[1cm]{C}。

		\fourchoices[answer=C]
		{内能}
		{分子平均动能}
		{单位体积内的分子数}
		{分子间相互作用力}
	\end{enumerate}

	\item
	为点燃硝化棉，需快速压下活塞，主要原因是时间较短可使封闭气体经筒壁 \tkanswer[1cm]{B}。

	\fourchoices[answer=B]
	{吸收的热量较少}
	{释放的热量较少}
	{释放的热量较多}
	{吸收的热量较多}

	\item
	引火仪的筒壁厚度对实验效果有影响。为测量筒壁厚度，将一束水平激光射入竖直放置的圆筒，调节入射角度，使光线恰好在筒壁内表面发生全反射，光路如图所示（俯视）。测得光线在筒壁外表面的入射角为 $\theta$，筒的外半径为 $R$，则筒壁厚度为 \tkanswer[2.2cm]{$R(1-\sin\theta)$}。

	\onepicture[h=4.5cm, align=right]{\includesvg{figs/04c}}
\end{enumerate}
%% answer:
\jdanswer{
\begin{enumerate}
	\item
	（1）AD     （2）C

	\item
	B

	\item
	$R(1-\sin\theta)$
\end{enumerate}
}
%% memo:
\memoanswer{
1．（1）汽缸内气体做等温压缩，在 $T-p$ 图像中等温线应是一根水平直线，选项 A 正确 B 错误；在 $V-p$ 图像中应是一根双曲线，选项 C 错误 D 正确。

正确选项为 AD。

（2）温度是物体分子热运动平均动能的量度，此过程是等温变化，分子平均动能不变。选项 B 错误。

一定质量理想气体的内能仅取决于它的温度，因此气体的内能不变。选项 A 错误。

理想气体不计分子间相互作用力，也就不计分子势能。选项 D 错误。

封闭气体的分子数保持不变，而体积减小，则单位体积内的分子数增大。选项 C 正确。

正确选项为 C。

2．快速压下活塞压缩气缸内的空气的过程可近似视为绝热变化，根据热力学第一定律，外力做功使气体内能增加，温度升高，而释放到外界的热量较少，硝化棉会更容易被点燃。

正确选项为 B。

3．光路如图所示，设折射角为 $\alpha$，根据折射率定义式可得 $n=\frac{\sin\theta}{\sin\alpha}$。

\onepicture[w=6.0cm]{\includesvg{figs/04_an}}

设全反射的临界角为 $C$，可得 $\sin C=\frac{1}{n}$。

设圆筒内表面半径为 $r$，在入射点、全反射点、圆心点构成的三角形中，根据正弦定理可得

\[\frac{r}{\sin\alpha}=\frac{R}{\sin(\pi-C)}=\frac{R}{\sin C}\]

解得 $r=\frac{R\sin\alpha}{\sin C}=R\frac{\frac{\sin\theta}{n}}{\frac{1}{n}}=R\sin\theta$。

则筒壁厚度为 $R-r=R(1-\sin\theta)$。
}

\item
%% number: 5
%% paperName: 2026年上海等级考第 5 题 20 分
%% typeId: 9
%% type: 综合题
%% chapter: 运动和力
%% point: 牛顿定律的应用
%% method: 无
%% score: 20
%% degree: 450
%% duplicateId: 0
%% body:
\textbf{松鼠}

一只松鼠在林间跳跃攀爬，其质量 $m=0.30 \Ukg$。

\onepicture[h=5.0cm, align=right]{\includesvg{figs/05a}}
\begin{enumerate}
	\item
	如图，该松鼠从一树枝的树梢沿水平方向跳出，跳到对面稍低处的另一水平树枝上。若两树枝在同一竖直平面内，且两树梢间的水平距离为 $L$、竖直距离为 $H$，则松鼠的跳出速度大小至少为 \tkanswer[1.8cm]{$L\sqrt{\frac{g}{2H}}$}。（重力加速度大小为 $g$，松鼠视为质点，忽略空气阻力）

	\item
	如图，该松鼠沿一树枝由静止开始向上做匀加速直线运动，速度方向与水平成 $\theta=60^{\circ}$ 的夹角，爬行 $2.0 \Um$ 用时 $2.0 \Us$。（忽略空气阻力，结果保留 2 位有效数字）

	\onepicture[h=2.8cm, align=right]{\includesvg{figs/05b}}

	\begin{enumerate}
		\item
		此过程中该松鼠爬行的加速度大小 $a=$ \tkanswer[1.2cm]{$1.0$} $\Umsq$。

		\item
		求此过程中树枝对该松鼠的作用力大小 $F$。（计算）
	\end{enumerate}

	\item
	小艾为研究松鼠从高处落地不受伤的原因，用高速摄像机拍摄了该松鼠竖直下落时的姿态，发现它在空中会伸开四肢和大尾巴，她测得该松鼠从高处由静止下落过程中的速度 $v$ 随下落距离 $h$ 变化关系如图所示。现研究该松鼠从 $10 \Um$ 高处由静止下落的情况。

	\onepicture[w=6.0cm, align=right]{\includesvg{figs/05c}}

	\begin{enumerate}
		\item
		该松鼠下落 $10 \Um$ 过程中，空气阻力对其所做的功为 \tkanswer[1.4cm]{$-17$} $\UJ$。（结果保留 2 位有效数字）

		\onepicture[h=2.8cm, align=right]{\includesvg{figs/05d}}

		\item
		小艾研究视频发现该松鼠接触松软草地后经 $t=90 \UmsT$ 速度减为零。若它可承受其所受重力 15 倍的平均冲击力而不受伤，通过计算证明它能安全落地。（论证）
	\end{enumerate}
\end{enumerate}
%% answer:
\jdanswer{
\begin{enumerate}
	\item
	$L\sqrt{\frac{g}{2H}}$

	\item
	（1）$1.0$     （2）$F=3.2 \UN$

	\item
	（1）$-17$     （2）取向上为正方向，由动量定理有 $(F-mg)t=0-(-mv)$，得 $F=mg+\frac{mv}{t}=9.8m+\frac{9m}{90\times10^{-3}}=109.8m<15mg=147m$，松鼠可以安全落地。
\end{enumerate}
}
%% memo:
\memoanswer{
1．松鼠做平抛运动，水平方向满足 $L=v_{0}t$，竖直方向满足 $H=\frac{1}{2}gt^{2}$，联立解得 $v_{0}=L\sqrt{\frac{g}{2H}}$。

2．（1）由 $x=\frac{1}{2}at^{2}$ 解得 $a=\frac{2x}{t^{2}}=\frac{2\times2}{2^{2}} \Umsq=1.0 \Umsq$。

（2）解法一：合成法

受力分析如图所示。

松鼠所受重力 $G=mg=0.3\times9.8 \UN=2.94 \UN$，所受合力 $F_{\text{合}}=ma=0.3 \UN$。

由余弦定理可得：

\[F=\sqrt{G^{2}+F_{\text{合}}^{2}-2GF_{\text{合}}\cos150^{\circ}}=\sqrt{2.94^{2}+0.3^{2}+2\times2.94\times0.3\times\frac{\sqrt{3}}{2}} \UN=3.2 \UN\]

\onepicture[h=4.0cm]{\includesvg{figs/05_an1}}

解法二：正交分解法

沿树枝方向和垂直树枝方向建立直角坐标系，受力分析如图所示。由正交分解法有：

$F\sin\alpha-G\cos60^{\circ}=0$，$F\cos\alpha-G\sin60^{\circ}=ma$。

又 $\sin^{2}\alpha+\cos^{2}\alpha=1$。

联立解得作用力大小为

\[F=m\sqrt{a^{2}+g^{2}+2ag\sin60^{\circ}}=0.3\times\sqrt{1.0^{2}+9.8^{2}+2\times1.0\times9.8\times\frac{\sqrt{3}}{2}} \UN=3.2 \UN\]

\onepicture[h=4.0cm]{\includesvg{figs/05_an2}}

3．（1）由图可知松鼠下落 $10 \Um$ 时的速度为 $9 \Ums$，由动能定理可得：

\[mgh+W_{f}=\frac{1}{2}mv^{2}-0\]

解得 $W_{f}=\frac{1}{2}mv^{2}-mgh=(0.5\times0.3\times9^{2}-0.3\times9.8\times10) \UJ=-17.25 \UJ$。

（2）取向上为正方向，由动量定理有 $(F-mg)t=0-(-mv)$，

得 $F=mg+\frac{mv}{t}=9.8m+m\frac{9}{90\times10^{-3}}=109.8m<15mg=147m$。

松鼠可以安全落地。
}

\item
%% number: 6
%% paperName: 2026年上海等级考第 6 题 21 分
%% typeId: 9
%% type: 综合题
%% chapter: 电磁感应
%% point: 电磁感应中的匀变速
%% method: 无
%% score: 21
%% degree: 350
%% duplicateId: 0
%% body:
\textbf{磁电式轮速传感器}

某磁电式轮速传感器基本结构如图所示。齿圈上形状相同的铁齿分布均匀。铁齿在磁铁附近被磁化，齿圈转动使穿过线圈的磁通量发生变化，从而通过电压传感器信号获得齿圈转速。

\onepicture[w=6.5cm, align=right]{\includesvg{figs/06a}}
\begin{enumerate}
	\item
	如图，齿圈从图示位置转过一个齿距的过程中，穿过线圈的磁通量先减小后增大。设 $a$、$b$ 两点电势分别为 $\varphi_{a}$、$\varphi_{b}$，则此过程中 \tkanswer[1cm]{D}。

	\onepicture[w=6.5cm, align=right]{\includesvg{figs/06b}}

	\fourchoices[answer=D]
	{$\varphi_{a}$ 始终高于 $\varphi_{b}$}
	{先 $\varphi_{a}$ 高于 $\varphi_{b}$，后 $\varphi_{a}$ 低于 $\varphi_{b}$}
	{$\varphi_{a}$ 始终低于 $\varphi_{b}$}
	{先 $\varphi_{a}$ 低于 $\varphi_{b}$，后 $\varphi_{a}$ 高于 $\varphi_{b}$}

	\item
	该传感器用于测量车速时，齿圈与汽车车轮以相同角速度转动。汽车匀速行驶。

	\begin{enumerate}
		\item
		线圈两端输出电压 $u$ 随时间 $t$ 的变化周期为 $2t_{0}$，如图（a）所示。已知车轮直径为 $D$，齿圈齿数为 $N$，则车速为 \tkanswer[2cm]{$\frac{\pi D}{2Nt_{0}}$}。（车轮不打滑）

		\onepicture[w=5.5cm, align=right]{\includesvg{figs/06c}}

		\item
		传感器测得的正弦交流电电压最大值为 $U_{m}$，经半波整流后的波形如图（b）所示，整流后电压的有效值为 \tkanswer[1.4cm]{$\frac{U_{m}}{2}$}。

		\onepicture[w=5.5cm, align=right]{\includesvg{figs/06d}}
	\end{enumerate}

	\item
	将齿圈换成如图（a）所示的金属笼，电阻不计的两铜质细圆环正对放置，与环面垂直的三根金属棒 $a_{1}b_{1}$、$a_{2}b_{2}$ 和 $a_{3}b_{3}$ 均匀分布在两环间且与环接触良好。在该笼绕 $OO'$ 转动过程中，金属棒依次整根穿过方向与棒垂直的有界匀强磁场，磁场范围很窄，棒在磁场中的速度可视为始终与磁场垂直，如图（b）所示。（虚线外磁场忽略不计，车轮不打滑）

	\onepicture[w=9.0cm, align=right]{\includesvg{figs/06e}}

	已知：金属棒长度均为 $L$、阻值均为 $r$，有界磁场的磁感应强度大小为 $B$。

	\begin{enumerate}
		\item
		当 $a_{1}b_{1}$ 切割磁感线时，两环间电压为 $U$，求此时

		\begin{steps}
			\item
			$a_{1}b_{1}$ 产生的感应电动势 $E$。（计算）

			\item
			$a_{1}b_{1}$ 所受安培力大小 $F$。（计算）
		\end{steps}

		\item
		笼与汽车车轮以相同角速度转动，两圆环半径均为 $R$，车轮半径是圆环半径的 $n$ 倍。某段时间内汽车做匀加速直线运动，测得两环间相继两次脉冲电压为 $U_{1}$ 和 $U_{2}$。

		\begin{steps}
			\item
			下一个脉冲电压 $U_{3}$ 为 \tkanswer[1cm]{D}。

			\fourchoices[answer=D]
			{$2U_{2}-U_{1}$}
			{$\frac{U_{2}^{2}}{U_{1}}$}
			{$\sqrt{\frac{U_{1}^{2}+U_{2}^{2}}{2}}$}
			{$\sqrt{2U_{2}^{2}-U_{1}^{2}}$}

			\item
			求汽车的加速度大小 $a$。（计算）
		\end{steps}
	\end{enumerate}
\end{enumerate}
%% answer:
\jdanswer{
\begin{enumerate}
	\item
	D

	\item
	（1）$\frac{\pi D}{2Nt_{0}}$     （2）$\frac{U_{m}}{2}$

	\item
	（1）① $E=3U$；② $F_{\text{安}}=\frac{2BUL}{r}$     （2）① D；② $a=\frac{27n(U_{2}^{2}-U_{1}^{2})}{4\pi B^{2}L^{2}R}$
\end{enumerate}
}
%% memo:
\memoanswer{
1．磁铁内磁感线方向由 S 极指向 N 极，原磁场方向向左，当磁通量减小时，根据楞次定律，感应电流的磁场方向也向左，根据右手螺旋定则判断出线圈中感应电流从 $a$ 流向 $b$，线圈作为一个电源，内部的电流是从低电势流向高电势，因此 $\varphi_{a}<\varphi_{b}$；同理，磁通量增大时，感应电流的方向从 $b$ 流向 $a$，$\varphi_{a}>\varphi_{b}$。

正确选项为 D。

2．（1）由图可知每隔 $2t_{0}$ 时间转过一个齿，则车轮转动的周期 $T=N\cdot2t_{0}$。由于车轮不打滑，因此车速即轮边缘的线速度，有

\[v=\frac{2\pi}{T}\cdot\frac{D}{2}=\frac{\pi D}{2Nt_{0}}\]

（2）经过半波整流的交流电在 $2t_{0}$ 时间内产生的热量为 $\frac{\left(\frac{U_{m}}{\sqrt{2}}\right)^{2}}{R}t_{0}$；在相同的电阻 $R$ 上加上直流电压 $U$，经过一个交变电流周期 $2t_{0}$ 的时间，产生的热量为 $\frac{U^{2}}{R}2t_{0}$，若满足热量相等：

\[\frac{\left(\frac{U_{m}}{\sqrt{2}}\right)^{2}}{R}t_{0}=\frac{U^{2}}{R}2t_{0}\]

得 $U=\frac{U_{m}}{2}$。

根据交流电有效值的定义，$U$ 即为交变电流的电压有效值。

3．（1）①金属棒 $a_{1}b_{1}$ 切割磁感线，作为电源，内电阻即为 $r$；另两根金属棒并联作为外电阻，外电阻 $R=\frac{r}{2}$。由闭合电路欧姆定律，有：

\[U=\frac{R}{R+r}E=\frac{1}{3}E\]

解得 $E=3U$。

②根据安培力公式可得 $F_{\text{安}}=BIL=B\cdot\frac{E}{R+r}\cdot L=B\cdot\frac{3U}{0.5r+r}\cdot L=\frac{2BUL}{r}$。

（2）①相继两次脉冲时间内汽车走过的距离相同，都是三分之一圆弧，又汽车做匀加速直线运动，有：

$v_{2}^{2}-v_{1}^{2}=2ax$，$v_{3}^{2}-v_{2}^{2}=2ax$，

得 $v_{3}=\sqrt{2v_{2}^{2}-v_{1}^{2}}$。

而电压 $U\propto E=BLv\propto v_{\text{车}}$，所以电压也有类似的规律，即 $U_{3}=\sqrt{2U_{2}^{2}-U_{1}^{2}}$。

正确选项为 D。

②笼与车轮角速度相同，由 $v=\omega r$ 可知车速为棒切割速度的 $n$ 倍。

由（1）小题结论可得 $E_{1}=3U_{1}=\frac{BLv_{1}}{n}$，$E_{2}=3U_{2}=\frac{BLv_{2}}{n}$。

则两次脉冲时刻车的速度分别为 $v_{1}=\frac{3nU_{1}}{BL}$、$v_{2}=\frac{3nU_{2}}{BL}$。

在这段时间内车轮转过三分之一圆周，车沿地面的位移为 $x=\frac{1}{3}\cdot2\pi nR$。

得加速度大小

\[a=\frac{v_{2}^{2}-v_{1}^{2}}{2x}=\frac{\left(\frac{3nU_{2}}{BL}\right)^{2}-\left(\frac{3nU_{1}}{BL}\right)^{2}}{\frac{4\pi nR}{3}}=\frac{27n(U_{2}^{2}-U_{1}^{2})}{4\pi B^{2}L^{2}R}\]
}

\end{enumerate}

\end{document}
